% Generated by book/tools/notebook_to_tex.py. Do not edit by hand.

\chapter[Multi-label]{다중 라벨과 per-label BCE}

\markboth{Multi-label}{Multi-label}

\chaptermeta{06_sklearn_multilabel/06_sklearn_multilabel.ipynb}{https://colab.research.google.com/github/yoon-gu/neuqes-101/blob/master/06_sklearn_multilabel/06_sklearn_multilabel.ipynb}{softmax의 합=1 제약을 풀고 라벨별 sigmoid로 확장}



\index{Multi-label classification}

\index{OneVsRestClassifier}

\index{BCEWithLogitsLoss}

\index{hamming\_loss}

\index{micro F1}

\index{macro F1}

\index{다중 라벨 분류}

\index{멀티핫}

\index{라벨별 BCE}

\index{해밍 손실}

\index{마이크로 F1}

\index{매크로 F1}

\index{임계값 조정}

\index{측면 라벨}



\section{변화추적표}

\begin{adjustbox}{max width=\textwidth}
\begin{tabular}{@{}lllllll@{}}
\toprule
Ch & 모델 & 토크나이저 & 데이터 & Output Head & Activation & Loss \\
\midrule

1 & (TF-IDF) & \inlinecode{TfidfVectorizer()} & Yelp 5,000 & --- & --- & --- \\
2 & \inlinecode{LinearRegression()} & \inlinecode{TfidfVectorizer()} & Yelp (별점 1-5) & (1차원) & 없음 & \inlinecode{MSELoss} \\
3 & \inlinecode{LogisticRegression()} & \inlinecode{TfidfVectorizer()} & Yelp 이진화 & (1차원) & sigmoid & \inlinecode{BCEWithLogitsLoss} \\
4 & \inlinecode{LogisticRegression(multi\_class="multinomial")} & \inlinecode{TfidfVectorizer()} & Yelp 이진화 (3장과 동일) & (2차원) & softmax & \inlinecode{CrossEntropyLoss} \\
5 & \inlinecode{LogisticRegression(multi\_class="multinomial")} & \inlinecode{TfidfVectorizer()} & Yelp 5클래스 & (5차원) & softmax & \inlinecode{CrossEntropyLoss} \\
\textbf{6 ← 여기} & \inlinecode{OneVsRestClassifier(LogisticRegression())} & \inlinecode{TfidfVectorizer()} & Yelp + 측면 키워드 합성 & (5차원) & \textbf{sigmoid (각각 독립)} & \textbf{\inlinecode{BCEWithLogitsLoss} per-label} \\
\bottomrule
\end{tabular}
\end{adjustbox}

전체 19개 장의 표는 \href{https://github.com/yoon-gu/neuqes-101\#챕터별-변화추적표}{루트 README.md}를 참고하세요.



\section{변경점: 5장 대비}

가장 큰 변화는 \textbf{데이터 가정} 이고, 그게 모델 선택을 강제합니다.

\begin{adjustbox}{max width=\textwidth}
\begin{tabular}{@{}lll@{}}
\toprule
축 & 5장 (multi-class) & 6장 (multi-label) \\
\midrule

데이터 가정 & 클래스 \emph{상호배타} (한 샘플 = 한 라벨) & \textbf{라벨 \emph{독립} (한 샘플에 여러 라벨 가능)} \\
라벨 형식 & int 한 개 (0-4) & \textbf{multi-hot 벡터} (예: \texttt{{[}1,\ 0,\ 1,\ 0,\ 1{]}}) \\
모델 패러다임 & \textbf{multinomial 기본} + OvR 대안 (5장 후반에서 두 방식 비교) & \textbf{OvR 만} (multinomial은 데이터 가정과 충돌) \\
Activation & softmax 한 번 (합 = 1 강제) & \textbf{per-label sigmoid} (라벨끼리 독립) \\
Loss & \inlinecode{CrossEntropyLoss} & \textbf{per-label \inlinecode{BCEWithLogitsLoss} 평균} \\
OvR 사용 방식 & K개 sigmoid → \textbf{argmax로 한 라벨 선택} & K개 sigmoid \textbf{그대로} (argmax 없음, 각자 임계값 0.5와 비교) \\
데이터 & 별점 5클래스 & \textbf{Yelp + 측면 키워드 합성} (5개 측면) \\
토크나이저 & TF-IDF & TF-IDF (그대로) \\
\bottomrule
\end{tabular}
\end{adjustbox}

\subsection{왜 OvR이 multi-label의 자연스러운 선택인가}

5장에선 두 패러다임이 모두 가능했습니다.

\begin{itemize}
\tightlist
\item
  \textbf{multinomial}: softmax 한 번으로 K개 logit을 묶어 합=1 강제. "K개 중 정확히 하나"라는 데이터 가정과 정합.
\item
  \textbf{OvR (대안)}: K개 \emph{독립} sigmoid + argmax 후처리. 각 binary 모델은 독립이지만 마지막에 강제로 하나만 고르므로 결과는 상호배타.
\end{itemize}

6장의 multi-label은 이 가정 자체를 깹니다 --- 한 리뷰가 "음식 + 서비스 + 가격"을 동시에 다룰 수 있어요. 그러면:

\begin{itemize}
\tightlist
\item
  \textbf{multinomial은 부적합}: softmax가 합=1을 강제하므로 \textquotesingle food=0.9, service=0.85 동시 활성\textquotesingle{} 같은 분포를 \emph{표현할 수가 없습니다}. P(food)=0.9면 나머지 합이 0.1로 강제돼 동시 활성이 수학적으로 불가능.
\item
  \textbf{OvR은 자연스럽게 들어맞음}: K개 sigmoid가 \emph{각자} 0/1을 결정 → 어떤 조합이든 표현 가능. argmax 후처리 단계만 빼면 그대로 multi-label.
\end{itemize}

요약: 5장에서 \emph{대안} 이었던 OvR이 6장에서는 \emph{유일한 자연스러운 선택} 이 됩니다. 알고리즘(\inlinecode{OneVsRestClassifier})은 그대로, 사용 방식만 "argmax로 한 라벨 선택" → "K개 출력 그대로" 로 바뀝니다.



\section{\texorpdfstring{Loss 함수의 변화 --- \inlinecode{BCEWithLogitsLoss} per-label}{Loss 함수의 변화 --- BCEWithLogitsLoss per-label}}

K개 라벨에 대해 BCE를 각각 계산하고 평균을 냅니다.

\begin{equation}
\label{eq:ch06-01}
L = \frac{1}{N \cdot K}\sum_{i=1}^{N}\sum_{k=1}^{K}\bigl[-y_{ik}\log \hat p_{ik} - (1 - y_{ik})\log(1 - \hat p_{ik})\bigr]
\end{equation}

식~\eqref{eq:ch06-01}은 정답 라벨과 예측 확률 사이의 Binary Cross-Entropy를 정의합니다.

각 (샘플, 라벨) 쌍이 독립적으로 손실에 기여합니다 --- 한 라벨에서 틀렸다고 다른 라벨의 손실이 변하지 않습니다. CE의 클래스 경쟁 구조와 정반대.

\textbf{숫자로 감 잡기} (K=5, 정답 multi-hot \(\mathbf{y} = [1, 0, 1, 0, 1]\)):

\begin{adjustbox}{max width=\textwidth}
\begin{tabular}{@{}llll@{}}
\toprule
시나리오 & 예측 확률 \(\hat{\mathbf{p}}\) & 라벨별 손실 & 평균 BCE \\
\midrule

잘 맞춤 & \texttt{{[}0.9,\ 0.1,\ 0.8,\ 0.2,\ 0.6{]}} & 0.105 / 0.105 / 0.223 / 0.223 / 0.511 & \textbf{0.233} \\
균등 (baseline) & \texttt{{[}0.5,\ 0.5,\ 0.5,\ 0.5,\ 0.5{]}} & 0.693 × 5 & \textbf{0.693} \\
정반대로 자신감 & \texttt{{[}0.1,\ 0.9,\ 0.1,\ 0.9,\ 0.1{]}} & 2.303 × 5 & \textbf{2.303} \\
\bottomrule
\end{tabular}
\end{adjustbox}

baseline = \(\log 2 = 0.693\) --- 모든 라벨에 0.5를 줄 때 (BCE에서 K=2 분포의 균등 추측과 같은 값). 학습된 모델은 이 값보다 작아야 정상.

\Needspace{14\baselineskip}
\begin{lstlisting}
# PyTorch (12장 이후, multi-label)
criterion = nn.BCEWithLogitsLoss()
loss = criterion(logits, targets.float())
# logits: (N, K), targets: (N, K) multi-hot

# sklearn (이 장)
from sklearn.multiclass import OneVsRestClassifier
model = OneVsRestClassifier(LogisticRegression(max_iter=1000))
model.fit(X, Y_multilabel)
# Y_multilabel shape: (N, K) 0/1
\end{lstlisting}



\section{토크나이저 노트}

이 장의 토크나이저는 \textbf{1-5장와 동일한 \inlinecode{TfidfVectorizer}}. 입력 표현은 그대로고, 변화는 라벨 구조와 출력 헤드의 형태에 있습니다.

\begin{previewBox}{미리보기}
\textbf{다음 장(7장)} --- Phase 1 시작: 사전학습된 \textbf{DistilBERT WordPiece} 가 처음 등장합니다. TF-IDF의 단어 단위 어휘 학습과 어떻게 다른지 비교 시작.
\end{previewBox}



\Needspace{8\baselineskip}
\begin{lstlisting}
!pip install -q datasets scikit-learn pandas matplotlib
\end{lstlisting}

\begin{codeRead}
\textbf{1행}의 \inlinecode{!pip install -q datasets scikit-learn pan...}에서는 Colab 실행에 필요한 패키지를 설치합니다.
\end{codeRead}



\Needspace{26\baselineskip}
\begin{lstlisting}
import warnings
warnings.filterwarnings("ignore", category=FutureWarning)
warnings.filterwarnings(
    "ignore",
    category=DeprecationWarning,
)

import re
import numpy as np
import pandas as pd
import matplotlib.pyplot as plt
from datasets import load_dataset
from sklearn.feature_extraction.text import TfidfVectorizer
from sklearn.linear_model import LogisticRegression
from sklearn.multiclass import OneVsRestClassifier
from sklearn.model_selection import train_test_split
from sklearn.metrics import (
    accuracy_score, hamming_loss, f1_score, classification_report,
)

plt.rcParams["axes.unicode_minus"] = False

dataset = load_dataset("yelp_review_full")
SAMPLE_SIZE = 5000
ds = dataset["train"].shuffle(seed=42).select(range(SAMPLE_SIZE))
df = ds.to_pandas()
display(pd.DataFrame([{"metric": "sample_count", "value": len(df)}]))
\end{lstlisting}

\begin{codeRead}
\inlinecode{numpy, pandas, matplotlib, datasets, scikit-learn} 같은 주요 패키지를 불러옵니다.
\end{codeRead}

\noindent\textbf{출력.}
\begin{lstlisting}[style=bookoutput]
metric  value
0  sample_count   5000
\end{lstlisting}
\par\vspace{0.9em}



\section{실습 1: 측면 라벨 합성}

Yelp 데이터에는 multi-label 정답이 없으므로 \textbf{5개 측면(aspect)} 별 키워드 사전을 만들어 매칭합니다.

\begin{adjustbox}{max width=\textwidth}
\begin{tabular}{@{}lll@{}}
\toprule
측면 & 의미 & 키워드 예시 \\
\midrule

\inlinecode{food} & 음식의 맛/메뉴 & food, meal, dish, taste, delicious, ... \\
\inlinecode{service} & 서비스/응대 & service, staff, waiter, friendly, rude, ... \\
\inlinecode{price} & 가격/가성비 & price, cheap, expensive, value, worth, ... \\
\inlinecode{ambiance} & 분위기/인테리어 & atmosphere, decor, music, vibe, cozy, ... \\
\inlinecode{location} & 위치/주차 & location, parking, area, neighborhood, ... \\
\bottomrule
\end{tabular}
\end{adjustbox}

각 리뷰 텍스트에 대해 \emph{어떤 키워드라도} 등장하면 해당 측면을 1로 활성화 --- 5차원 multi-hot 벡터가 됩니다. \textbf{이 합성의 한계} 는 장 끝에서 솔직히 짚습니다.



\Needspace{26\baselineskip}
\begin{lstlisting}
ASPECT_KEYWORDS = {
    "food": ["food", "meal", "dish", "taste", "delicious", "flavor", "menu",
             "cuisine", "tasty", "yummy", "spicy", "sweet", "salty", "fresh"],
    "service": ["service", "staff", "waiter", "waitress", "server", "friendly",
                "rude", "attentive", "host", "helpful", "polite", "manager"],
    "price": ["price", "cheap", "expensive", "value", "worth", "cost",
              "money", "afford", "overpriced", "pricey", "deal", "bargain"],
    "ambiance": ["atmosphere", "ambiance", "decor", "music", "vibe", "cozy",
                 "noisy", "quiet", "lighting", "interior", "comfortable", "loud"],
    "location": ["location", "parking", "area", "neighborhood", "access",
                 "downtown", "convenient", "spot"],
}
ASPECTS = list(ASPECT_KEYWORDS.keys())
K = len(ASPECTS)

def extract_aspects(text: str) -> list[int]:
    tokens = set(re.findall(r"[a-z]+", text.lower()))
    return [
        int(any(kw in tokens for kw in keywords))
        for keywords in ASPECT_KEYWORDS.values()
    ]

# 5,000건에 적용
df["aspects"] = df["text"].apply(extract_aspects)
Y = np.array(df["aspects"].tolist())   # (N, 5) multi-hot

label_shape = pd.DataFrame([{"object": "Y", "shape": Y.shape}])
label_preview = pd.DataFrame(Y[:3], columns=ASPECTS)

display(label_shape)
display(label_preview)
\end{lstlisting}

\begin{codeRead}
\textbf{1--12행}의 \inlinecode{ASPECT\_KEYWORDS = \{ "food": ["food", "mea...}에서는 이후 분석에서 사용할 값을 준비합니다. \textbf{13행}의 \inlinecode{ASPECTS = list(ASPECT\_KEYWORDS.keys())}에서는 이후 분석에서 사용할 값을 준비합니다. \textbf{14행}의 \inlinecode{K = len(ASPECTS)}에서는 이후 분석에서 사용할 값을 준비합니다. \textbf{16행}의 \inlinecode{def extract\_aspects(text: str) -> list[int]:}에서는 앞 단계에서 만든 값을 바탕으로 다음 계산을 수행합니다. \textbf{17행}의 \inlinecode{tokens = set(re.findall(r"[a-z]+", text.l...}에서는 이후 분석에서 사용할 값을 준비합니다. \textbf{18--21행}의 \inlinecode{return [ int(any(kw in tokens for kw in k...}에서는 앞 단계에서 만든 값을 바탕으로 다음 계산을 수행합니다.
\end{codeRead}

\noindent\textbf{출력.}
\begin{lstlisting}[style=bookoutput]
object      shape
0      Y  (5000, 5)
   food  service  price  ambiance  location
0     0        1      0         0         1
1     0        0      0         0         1
2     0        0      0         1         0
\end{lstlisting}
\par\vspace{0.9em}



\Needspace{21\baselineskip}
\begin{lstlisting}
aspect_activation = pd.DataFrame({
    "aspect": ASPECTS,
    "activation_rate": [Y[:, k].mean() for k in range(K)],
    "active_count": [int(Y[:, k].sum()) for k in range(K)],
})

n_labels_per_sample = Y.sum(axis=1)
label_count_distribution = pd.DataFrame({
    "n_active_labels": range(K + 1),
    "sample_count": [(n_labels_per_sample == n).sum() for n in range(K + 1)],
})
label_count_distribution["sample_rate"] = label_count_distribution["sample_count"] / len(Y)
summary = pd.DataFrame([{"metric": "mean_active_labels_per_sample", "value": n_labels_per_sample.mean()}])

display(aspect_activation)
display(summary)
display(label_count_distribution)
\end{lstlisting}

\begin{codeRead}
\textbf{1--5행}의 \inlinecode{aspect\_activation = pd.DataFrame(\{ "aspec...}에서는 결과를 표로 보기 좋게 정리합니다. \textbf{7행}의 \inlinecode{n\_labels\_per\_sample = Y.sum(axis=1)}에서는 이후 분석에서 사용할 값을 준비합니다. \textbf{8--11행}의 \inlinecode{label\_count\_distribution = pd.DataFrame(\{...}에서는 결과를 표로 보기 좋게 정리합니다. \textbf{12행}의 \inlinecode{label\_count\_distribution["sample\_rate"] =...}에서는 이후 분석에서 사용할 값을 준비합니다. \textbf{13행}의 \inlinecode{summary = pd.DataFrame([\{"metric": "mean\_...}에서는 결과를 표로 보기 좋게 정리합니다.
\end{codeRead}

\noindent\textbf{출력.}
\begin{lstlisting}[style=bookoutput]
aspect  activation_rate  active_count
0      food           0.5556          2778
1   service           0.4960          2480
2     price           0.2944          1472
3  ambiance           0.1810           905
4  location           0.2190          1095
                          metric  value
0  mean_active_labels_per_sample  1.746
   n_active_labels  sample_count  sample_rate
0                0           741       0.1482
1                1          1464       0.2928
2                2          1506       0.3012
3                3           957       0.1914
4                4           277       0.0554
5                5            55       0.0110
\end{lstlisting}
\par\vspace{0.9em}



\Needspace{18\baselineskip}
\begin{lstlisting}
X_text_train, X_text_test, Y_train, Y_test = train_test_split(
    df["text"], Y,
    test_size=0.2, random_state=42,
)

tfidf = TfidfVectorizer(max_features=10000)
X_train = tfidf.fit_transform(X_text_train)
X_test = tfidf.transform(X_text_test)

split_summary = pd.DataFrame([
    {"split": "train", "X_shape": X_train.shape, "Y_shape": Y_train.shape},
    {"split": "test", "X_shape": X_test.shape, "Y_shape": Y_test.shape},
])
display(split_summary)
\end{lstlisting}

\begin{codeRead}
\textbf{1--4행}의 \inlinecode{X\_text\_train, X\_text\_test, Y\_train, Y\_tes...}에서는 학습용 데이터와 평가용 데이터를 분리합니다. \textbf{6행}의 \inlinecode{tfidf = TfidfVectorizer(max\_features=10000)}에서는 TF-IDF 기반 벡터라이저를 만듭니다. \textbf{7행}의 \inlinecode{X\_train = tfidf.fit\_transform(X\_text\_train)}에서는 텍스트나 라벨을 모델이 처리할 수 있는 수치 표현으로 변환합니다. \textbf{8행}의 \inlinecode{X\_test = tfidf.transform(X\_text\_test)}에서는 텍스트나 라벨을 모델이 처리할 수 있는 수치 표현으로 변환합니다. \textbf{10--13행}의 \inlinecode{split\_summary = pd.DataFrame([ \{"split":...}에서는 결과를 표로 보기 좋게 정리합니다.
\end{codeRead}

\noindent\textbf{출력.}
\begin{lstlisting}[style=bookoutput]
split        X_shape    Y_shape
0  train  (4000, 10000)  (4000, 5)
1   test  (1000, 10000)  (1000, 5)
\end{lstlisting}
\par\vspace{0.9em}



\section{\texorpdfstring{실습 2: \inlinecode{OneVsRestClassifier} (이번엔 argmax 없이)}{실습 2: OneVsRestClassifier (이번엔 argmax 없이)}}

5장에서 multi-class의 \emph{대안} 으로 본 \inlinecode{OneVsRestClassifier}가 이번엔 \emph{기본 도구} 입니다. 코드 차이는 단 하나 --- \textbf{\inlinecode{Y\_train} shape이 \inlinecode{(N,)} → \inlinecode{(N,\ K)}} 로 바뀌는 것. sklearn이 multi-hot Y를 보고 multi-label 모드로 자동 전환해 K개 binary 분류기가 독립적으로 학습됩니다.

가장 중요한 사용 방식 변화는 \textbf{argmax가 사라진다는 것}.

\begin{itemize}
\tightlist
\item
  \textbf{5장 OvR (multi-class)}: 5개 sigmoid 출력 중 가장 큰 것 \emph{하나만} 정답으로 골랐음.
\item
  \textbf{6장 OvR (multi-label)}: 5개 sigmoid 출력을 \emph{각자} 임계값 0.5와 비교해 0/1 결정. 동시 활성 가능.
\end{itemize}

\inlinecode{predict\_proba}도 차이가 있습니다 --- 5장에서는 sklearn이 후처리로 합=1로 정규화해줬지만, multi-label 모드에서는 \emph{정규화하지 않고} 각 라벨의 P(label\_k = 1)을 그대로 반환합니다 (각 라벨이 독립이니까).



\Needspace{14\baselineskip}
\begin{lstlisting}
# 먼저 wrapper 없이 단순히 LogisticRegression() 에 
# multi-hot Y를 넣어보기
bare_model = LogisticRegression(max_iter=1000)
try:
    bare_model.fit(X_train, Y_train)
    bare_result = {"status": "success", "detail": str(bare_model.coef_.shape)}
except ValueError as e:
    bare_result = {"status": "ValueError", "detail": str(e)}

display(pd.DataFrame([bare_result]))
\end{lstlisting}

\begin{codeRead}
\textbf{3행}의 \inlinecode{bare\_model = LogisticRegression(max\_iter=...}에서는 이번 실습에서 관찰할 모델 객체를 정의합니다. \textbf{4행}의 \inlinecode{try:}에서는 앞 단계에서 만든 값을 바탕으로 다음 계산을 수행합니다. \textbf{5행}의 \inlinecode{bare\_model.fit(X\_train, Y\_train)}에서는 학습 데이터로 모델 또는 변환기를 적합합니다. \textbf{6행}의 \inlinecode{bare\_result = \{"status": "success", "deta...}에서는 이후 분석에서 사용할 값을 준비합니다. \textbf{7행}의 \inlinecode{except ValueError as e:}에서는 앞 단계에서 만든 값을 바탕으로 다음 계산을 수행합니다. \textbf{8행}의 \inlinecode{bare\_result = \{"status": "ValueError", "d...}에서는 이후 분석에서 사용할 값을 준비합니다.
\end{codeRead}

\noindent\textbf{출력.}
\begin{lstlisting}[style=bookoutput]
status                                             detail
0  ValueError  y should be a 1d array, got an array of
  shape ...
\end{lstlisting}
\par\vspace{0.9em}



\textbf{왜 실패했나?} sklearn \inlinecode{LogisticRegression} 은 \emph{1D Y} (각 샘플당 한 클래스 인덱스)만 받습니다. 우리의 \inlinecode{Y\_train.shape\ ==\ (N,\ 5)} 는 "한 샘플에 여러 라벨"이라는 의미인데 \emph{형식 자체} 가 호환되지 않아요. fit 호출이 즉시 중단될 수 있습니다.

\inlinecode{OneVsRestClassifier} 의 역할은 단순합니다.

\begin{enumerate}
\def\labelenumi{\arabic{enumi}.}
\tightlist
\item
  2D Y를 K개의 1D 컬럼으로 \textbf{쪼개고},
\item
  각 컬럼마다 \inlinecode{LogisticRegression} 을 \textbf{별도로} 학습 (총 K개 모델),
\item
  결과를 \inlinecode{.estimators\_} 리스트에 보관해 \inlinecode{predict} 시 모두 적용.
\end{enumerate}

알고리즘은 동일한 LogReg지만 \textbf{fit 시점의 Y 형식 처리} 가 결정적 차이입니다 --- bare 호출은 실패하고, wrapper 호출은 K개 모델로 분할 학습됩니다.



\Needspace{26\baselineskip}
\begin{lstlisting}
# 위 실패와 대비: wrapper 한 줄로 K개 LogReg가 자동 분할 
# 학습됨
model_ml = OneVsRestClassifier(LogisticRegression(max_iter=1000))
model_ml.fit(X_train, Y_train)

first_estimator = model_ml.estimators_[0]
first_coef_shape = getattr(
    getattr(first_estimator, "coef_", None),
    "shape",
    None,
)
estimator_summary = pd.DataFrame([
    {"metric": "n_binary_estimators", "value": len(model_ml.estimators_)},
    {"metric": "estimator_0_type", "value": type(first_estimator).__name__},
    {"metric": "estimator_0_coef_shape", "value": first_coef_shape},
])
label_estimator_summary = pd.DataFrame([
    {
        "estimator": k,
        "aspect": aspect,
        "positive_count": int(Y_train[:, k].sum()),
        "positive_rate": Y_train[:, k].mean(),
    }
    for k, aspect in enumerate(ASPECTS)
])

# 예측 + 확률
Y_pred = model_ml.predict(X_test)
# (N, K) multi-hot 0/1 (threshold 0.5 자동 적용)
proba_ml = model_ml.predict_proba(X_test) # (N, K) per-label probability (정규화 X)

prediction_shape = pd.DataFrame([
    {"object": "Y_pred", "shape": Y_pred.shape},
    {"object": "proba_ml", "shape": proba_ml.shape},
])
proba_preview = pd.DataFrame(
    proba_ml[:3],
    columns=ASPECTS).round(4,
)

display(estimator_summary)
display(label_estimator_summary)
display(prediction_shape)
display(proba_preview)
\end{lstlisting}

\begin{codeRead}
\textbf{3행}의 \inlinecode{model\_ml = OneVsRestClassifier(LogisticRe...}에서는 이번 실습에서 관찰할 모델 객체를 정의합니다. \textbf{4행}의 \inlinecode{model\_ml.fit(X\_train, Y\_train)}에서는 학습 데이터로 모델 또는 변환기를 적합합니다. \textbf{6행}의 \inlinecode{first\_estimator = model\_ml.estimators\_[0]}에서는 이후 분석에서 사용할 값을 준비합니다. \textbf{7--11행}의 \inlinecode{first\_coef\_shape = getattr( getattr(first...}에서는 이후 분석에서 사용할 값을 준비합니다. \textbf{12--16행}의 \inlinecode{estimator\_summary = pd.DataFrame([ \{"metr...}에서는 결과를 표로 보기 좋게 정리합니다. \textbf{17--25행}의 \inlinecode{label\_estimator\_summary = pd.DataFrame([...}에서는 결과를 표로 보기 좋게 정리합니다.
\end{codeRead}

\noindent\textbf{출력.}
\begin{lstlisting}[style=bookoutput]
metric               value
0     n_binary_estimators                   5
1        estimator_0_type  LogisticRegression
2  estimator_0_coef_shape          (1, 10000)
   estimator    aspect  positive_count  positive_rate
0          0      food            2214        0.55350
1          1   service            1961        0.49025
2          2     price            1196        0.29900
3          3  ambiance             720        0.18000
4          4  location             892        0.22300
     object      shape
0    Y_pred  (1000, 5)
1  proba_ml  (1000, 5)
     food  service   price  ambiance  location
0  0.3002   0.3082  0.7970    0.1701    0.1122
1  0.3607   0.3612  0.4585    0.1855    0.1662
2  0.2121   0.2894  0.1212    0.1387    0.0692
\end{lstlisting}
\par\vspace{0.9em}



\section{Loss 한 단계 더: 학습된 모델의 실제 예측으로 BCE 분해}

방금 fit한 \inlinecode{model\_ml}이 한 샘플에 대해 어떤 손실을 만들어내는지 직접 분해합니다. 변경점 표에서 본 "\textbf{per-label BCE 평균}" 이 단순한 수식이 아니라 \textbf{실제 5개 숫자의 산수} 라는 걸 확인합니다 --- 그리고 그 위에서 "왜 multinomial CE는 여기 못 쓰나"를 진짜 값으로 짚습니다.



\Needspace{26\baselineskip}
\begin{lstlisting}
# 여러 라벨이 활성된 test 샘플 하나 고르기 (분해가 잘 
# 보이도록)
multi_active = np.where(Y_test.sum(axis=1) >= 3)[0]
sample_idx = int(multi_active[0]) if len(multi_active) > 0 else 0

y_true = Y_test[sample_idx]
p_pred = proba_ml[sample_idx]
text = X_text_test.iloc[sample_idx]

rows = []
total_loss = 0.0
for k, aspect in enumerate(ASPECTS):
    y_k, p_k = int(y_true[k]), float(p_pred[k])
    if y_k == 1:
        loss_k = -np.log(max(p_k, 1e-12))
        formula = f"-log({p_k:.4f})"
    else:
        loss_k = -np.log(max(1 - p_k, 1e-12))
        formula = f"-log(1-{p_k:.4f})"
    total_loss += loss_k
    rows.append({
        "aspect": aspect,
        "target": y_k,
        "predicted_probability": p_k,
        "loss_term": formula,
        "bce": loss_k,
    })

review_preview = pd.DataFrame([{"review_preview_200_chars": text[:200] + "..."}])
sample_summary = pd.DataFrame([{"active_label_count": int(y_true.sum())}])
bce_breakdown = pd.DataFrame(rows)
bce_summary = pd.DataFrame([
    {"metric": "sum_bce", "value": total_loss},
    {"metric": "mean_bce", "value": total_loss / K},
])

display(review_preview)
display(sample_summary)
display(bce_breakdown.round(4))
display(bce_summary.round(4))
\end{lstlisting}

\noindent\textbf{출력.}
\begin{lstlisting}[style=bookoutput]
review_preview_200_chars
0  Okay, so I just went to Vegas for the first ti...
   active_label_count
0                   3
     aspect  target  predicted_probability       loss_term
       bce
0      food       0                 0.4466  -log(1-0.4466)
  0.5917
1   service       1                 0.9027    -log(0.9027)
  0.1024
2     price       1                 0.7366    -log(0.7366)
  0.3057
3  ambiance       1                 0.5771    -log(0.5771)
  0.5497
4  location       0                 0.2620  -log(1-0.2620)
  0.3038
     metric   value
0   sum_bce  1.8533
1  mean_bce  0.3707
\end{lstlisting}
\par\vspace{0.9em}



\textbf{관찰}

\begin{itemize}
\tightlist
\item
  5개 라벨이 \emph{각자 독립적으로} 손실을 기여합니다 --- 한 라벨에서 잘 맞춰도 다른 라벨에서 못 맞추면 그 영향이 그대로 평균에 더해집니다.
\item
  정답이 1인 라벨은 \(-\log(p)\) --- 예측 확률이 1에 가까울수록 손실 0에 수렴.
\item
  정답이 0인 라벨은 \(-\log(1-p)\) --- 예측 확률이 0에 가까울수록 손실 0에 수렴.
\item
  같은 sigmoid 출력에 대해 정답이 0이냐 1이냐에 따라 \textbf{정반대 방향} 으로 페널티가 커집니다 (대칭 구조).
\end{itemize}

\subsection{같은 샘플을 multinomial CE로 풀려고 하면}

위 샘플의 정답은 multi-hot --- 여러 라벨이 동시에 1입니다. 만약 multinomial CE를 \emph{억지로} 적용하려면 다음 두 가지 \emph{임의 결정} 이 필요합니다.

\begin{enumerate}
\def\labelenumi{\arabic{enumi}.}
\tightlist
\item
  \textbf{5개 활성 라벨 중 \emph{하나만} 정답으로 골라야 함}: argmax? 첫 활성? 어느 기준이든 \emph{임의}.
\item
  \textbf{그러면 나머지 활성 라벨들은 \emph{틀린} 클래스로 학습됨}: 모델이 그 라벨에 강한 확률을 줄수록 손실이 \emph{커짐}.
\end{enumerate}

결과: 모델이 "동시 활성 패턴을 \emph{피하려고}" 학습됩니다 --- 실제 정답에서는 동시 활성이 정답인데도. 이건 \emph{데이터 가정과 정반대 방향} 으로 학습 신호가 작동하는 셈입니다.

per-label BCE는 위 표처럼 5개 손실을 \emph{독립적으로} 합산하므로 각 라벨이 자기 정답에만 책임을 집니다. \textbf{multi-label 데이터의 본래 구조와 정합한 유일한 선택} 인 이유가 이 산수에 있습니다.



\section{해부: multi-label 평가 지표}

multi-class에서 자주 쓰던 accuracy는 multi-label에서 의미가 미묘하게 다릅니다.

\begin{itemize}
\tightlist
\item
  \textbf{subset accuracy} (\inlinecode{accuracy\_score}): "K개 라벨이 \emph{전부} 일치한 샘플 비율" --- 라벨 하나만 틀려도 0점. 매우 엄격.
\item
  \textbf{hamming loss}: 평균 라벨별 오답 비율. 5개 중 1개 틀리면 0.2 기여. 가장 직관적.
\item
  \textbf{micro F1}: 모든 라벨의 TP/FP/FN를 한 풀에 모아 계산 --- 빈도 큰 라벨이 영향력 큼.
\item
  \textbf{macro F1}: 라벨별 F1을 단순 평균 --- 모든 라벨 동등 가중.
\end{itemize}



\Needspace{11\baselineskip}
\begin{lstlisting}
multi_label_metrics = pd.DataFrame([
    {"metric": "subset_accuracy", "value": accuracy_score(Y_test, Y_pred)},
    {"metric": "hamming_loss", "value": hamming_loss(Y_test, Y_pred)},
    {"metric": "micro_f1", "value": f1_score(Y_test, Y_pred, average="micro", zero_division=0)},
    {"metric": "macro_f1", "value": f1_score(Y_test, Y_pred, average="macro", zero_division=0)},
])
display(multi_label_metrics)
\end{lstlisting}

\begin{codeRead}
\textbf{1--6행}의 \inlinecode{multi\_label\_metrics = pd.DataFrame([ \{"me...}에서는 결과를 표로 보기 좋게 정리합니다.
\end{codeRead}

\noindent\textbf{출력.}
\begin{lstlisting}[style=bookoutput]
metric     value
0  subset_accuracy  0.493000
1     hamming_loss  0.138200
2         micro_f1  0.773665
3         macro_f1  0.646973
\end{lstlisting}
\par\vspace{0.9em}



\Needspace{13\baselineskip}
\begin{lstlisting}
report = classification_report(
    Y_test,
    Y_pred,
    target_names=ASPECTS,
    zero_division=0,
    output_dict=True,
)
report_df = pd.DataFrame(report).T
display(report_df.round(4))
\end{lstlisting}

\begin{codeRead}
\textbf{1--7행}의 \inlinecode{report = classification\_report( Y\_test, Y...}에서는 이후 분석에서 사용할 값을 준비합니다. \textbf{8행}의 \inlinecode{report\_df = pd.DataFrame(report).T}에서는 결과를 표로 보기 좋게 정리합니다.
\end{codeRead}

\noindent\textbf{출력.}
\begin{lstlisting}[style=bookoutput]
precision  recall  f1-score  support
food             0.8959  0.9007    0.8983    564.0
service          0.8936  0.8574    0.8751    519.0
price            0.9291  0.4275    0.5856    276.0
ambiance         1.0000  0.2757    0.4322    185.0
location         0.9365  0.2906    0.4436    203.0
micro avg        0.9043  0.6760    0.7737   1747.0
macro avg        0.9310  0.5504    0.6470   1747.0
weighted avg     0.9162  0.6760    0.7398   1747.0
samples avg      0.6832  0.5788    0.6082   1747.0
\end{lstlisting}
\par\vspace{0.9em}



\section{변형: 라벨별 임계값(threshold) 변경}

multi-label에서는 K개 sigmoid 출력에 대해 임계값 0.5로 자르는 게 기본입니다. 이 임계값을 옮기면 모든 라벨이 함께 반응합니다 --- 임계값을 낮추면 더 많은 라벨이 활성되어 recall은 오르고 precision은 내려갑니다.

(고급 트릭: 라벨마다 \emph{별도} 임계값을 정해 검증 F1을 최대화할 수도 있음. 여기서는 하나로 통일.)



\Needspace{17\baselineskip}
\begin{lstlisting}
thresholds = [0.2, 0.3, 0.4, 0.5, 0.6, 0.7]
rows = []
for t in thresholds:
    Y_pred_t = (proba_ml >= t).astype(int)
    rows.append({
        "threshold": t,
        "subset_acc": accuracy_score(Y_test, Y_pred_t),
        "hamming": hamming_loss(Y_test, Y_pred_t),
        "micro_F1": f1_score(Y_test, Y_pred_t, average="micro", zero_division=0),
        "macro_F1": f1_score(Y_test, Y_pred_t, average="macro", zero_division=0),
    })
df_t = pd.DataFrame(rows).round(4)
display(df_t)
\end{lstlisting}

\noindent\textbf{출력.}
\begin{lstlisting}[style=bookoutput]
threshold  subset_acc  hamming  micro_F1  macro_F1
0        0.2       0.233   0.2696    0.7178    0.6984
1        0.3       0.472   0.1468    0.8135    0.7877
2        0.4       0.563   0.1136    0.8335    0.7697
3        0.5       0.493   0.1382    0.7737    0.6470
4        0.6       0.416   0.1658    0.6984    0.5129
5        0.7       0.344   0.2028    0.5947    0.3919
\end{lstlisting}
\par\vspace{0.9em}



\section{합성의 한계 --- 솔직한 한계 짚기}

이 장의 학습 결과가 너무 좋아 보일 수 있습니다 (subset accuracy, micro F1 모두 매우 높음). 그 이유는 \textbf{모델이 \emph{키워드 매칭 규칙} 자체를 학습} 하기 때문입니다 --- 우리가 정한 사전을 다시 거꾸로 풀어내고 있을 뿐, 진짜 측면 추출 능력을 입증한 게 아닙니다.

실제 multi-label 문제에서 부딪히는 것들:

\begin{enumerate}
\def\labelenumi{\arabic{enumi}.}
\tightlist
\item
  \textbf{부정·반어 무시} --- \inlinecode{"this\ place\ is\ not\ noisy\ at\ all"} 의 \textquotesingle noisy\textquotesingle 를 ambiance 활성으로 잡는 게 우리 사전의 한계. 사람이 읽으면 ambiance가 \emph{아닌} 데도.
\item
  \textbf{사전 협소} --- \textquotesingle food\textquotesingle 에 \textquotesingle sushi\textquotesingle, \textquotesingle ramen\textquotesingle, \textquotesingle pasta\textquotesingle{} 같은 구체 음식명이 빠져 있으면 그 리뷰는 food=0이 되어 버림.
\item
  \textbf{정답에 노이즈가 있음} --- 우리 라벨 자체가 진짜 정답이 아닌 휴리스틱이라, 모델 성능을 이 정답에 비교하는 건 결국 "모델이 휴리스틱을 얼마나 따라 했나"를 잴 뿐.
\item
  \textbf{빈 라벨}: 모든 측면이 0인 샘플도 있음 (\texttt{\{n\_labels\_per\_sample\ ==\ 0).sum()} 건). 실제 multi-label 데이터에선 보통 최소 한 라벨은 보장.
\end{enumerate}

\textbf{그럼 왜 합성을 쓰나?} --- 학습 코드의 \emph{형태} 와 \emph{평가 지표 해석} 을 익히는 게 이 장의 목적이기 때문입니다. 12장 BERT multi-label에서 \textbf{같은 합성 라벨} 을 그대로 사용하므로 비교가 일관되게 됩니다. 진짜 multi-label 데이터(예: GoEmotions, Reuters)는 라벨이 사람 손으로 만들어져 있어 비용이 큽니다.



\section{이 장에 등장한 라이브러리}

\begin{adjustbox}{max width=\textwidth}
\begin{tabular}{@{}lll@{}}
\toprule
이름 & 한 줄 설명 & 다음 장에서 \\
\midrule

\inlinecode{sklearn.multiclass.OneVsRestClassifier} & K개 독립 binary 분류기를 묶고 multi-label 모드 자동 감지 & 12장 BERT multi-label에서 같은 패러다임을 BERT로 \\
\inlinecode{sklearn.metrics.hamming\_loss} & 라벨별 평균 오답 비율 & 12장에서도 평가 지표로 \\
\inlinecode{sklearn.metrics.f1\_score(average="micro"\ /\ "macro")} & multi-label F1 집계 방식 & --- \\
\inlinecode{sklearn.metrics.classification\_report} & 라벨별 precision/recall/F1 한 번에 & --- \\
\bottomrule
\end{tabular}
\end{adjustbox}



\section{체크포인트 질문}

\begin{enumerate}
\def\labelenumi{\arabic{enumi}.}
\tightlist
\item
  multi-class와 multi-label은 라벨 구조와 활성화 함수에서 어떻게 다른가요?
\item
  multi-label에서 subset accuracy가 너무 엄격한 지표가 되는 이유는?
\item
  multi-label baseline BCE는 왜 \(\log 2 = 0.693\) 인가요?
\item
  임계값을 0.5에서 0.3으로 낮추면 micro F1과 macro F1은 어느 방향으로 움직이는 게 일반적이고, 그 이유는 무엇인가요?
\end{enumerate}



\section{FAQ}

\begin{faqBox}{Q1. (이론) Multi-class와 Multi-label은 정확히 어떻게 다른가요?}

\begin{adjustbox}{max width=\textwidth}
\begin{tabular}{@{}lll@{}}
\toprule
항목 & Multi-class & Multi-label \\
\midrule

한 샘플의 라벨 수 & 정확히 1개 & 0개 이상 (K개 가능) \\
라벨 구조 & 정수 인덱스 (예: 3) & multi-hot 벡터 (예: {[}1, 0, 1, 0, 1{]}) \\
활성화 & softmax (합 = 1) & per-label sigmoid (라벨끼리 독립) \\
Loss & CrossEntropy & per-label BCE 평균 \\
가정 & 클래스 \emph{상호배타} & 라벨 \emph{독립} \\
예시 & 별점 1-5 중 하나, 뉴스 7카테고리 & 영화 장르(로맨스+코미디), 영화 태그 \\
\bottomrule
\end{tabular}
\end{adjustbox}

\end{faqBox}

\begin{faqBox}{Q2. (이론) micro F1과 macro F1 중 뭘 봐야 하나요?}

데이터 분포에 따라 다릅니다.

\begin{itemize}
\tightlist
\item
  \textbf{micro F1}: 라벨 빈도와 무관하게 \emph{전체 예측 풀} 의 정확도를 봄. 빈도 큰 라벨이 점수를 끌고감 --- 우리 데이터에서 food가 매우 자주 활성된다면 food 성능이 좋으면 micro F1도 좋아 보임.
\item
  \textbf{macro F1}: 라벨별 F1을 단순 평균. \emph{드문 라벨} 도 동등하게 평가 --- location 같은 빈도 작은 라벨에서 못하면 점수가 깎임.
\end{itemize}

\textbf{언제 무엇:}

\begin{itemize}
\tightlist
\item
  라벨 빈도가 비슷 + 모든 라벨 동등 중요 → 둘이 비슷, 단순히 macro 보면 됨.
\item
  빈도 차이가 크고 \emph{드문 라벨도 잘 잡고 싶다} → macro F1 (소수 클래스 보호).
\item
  전체 시스템의 평균적 정확도가 핵심 → micro F1.
\end{itemize}

실무에선 둘 다 보고 차이를 해석하는 게 표준입니다.

\end{faqBox}

\begin{faqBox}{Q3. (실무) 모든 라벨이 0인 샘플은 어떻게 처리하나요?}

세 가지 접근.

\begin{enumerate}
\def\labelenumi{\arabic{enumi}.}
\tightlist
\item
  \textbf{그대로 둔다}: BCE가 처리 가능. 각 라벨이 0이라는 신호를 학습. 가장 흔한 방식.
\item
  \textbf{버린다}: "라벨이 없으면 학습 신호가 없다"는 가정. 합성 데이터에서 빈 라벨 비율이 너무 크면 고려.
\item
  \textbf{"기타" 라벨 추가}: K+1번째 라벨을 만들어 "어느 측면도 안 맞음"을 명시적으로 표현. 데이터셋 설계 결정.
\end{enumerate}

\Needspace{8\baselineskip}
\begin{lstlisting}
# 옵션 (b): 빈 라벨 샘플 제거
mask_nonempty = Y.sum(axis=1) > 0
Y_clean = Y[mask_nonempty]
df_clean = df[mask_nonempty]
\end{lstlisting}

이 장은 (a) 그대로 두기 --- 합성 데이터에서 "어느 키워드도 없는 짧은 리뷰"는 자연스러운 부분이라.

\end{faqBox}

\begin{faqBox}{Q4. (실무) 키워드 매칭이 너무 단순한 것 같은데 실무에선 어떻게 하나요?}

세 가지 접근이 일반적.

\begin{enumerate}
\def\labelenumi{\arabic{enumi}.}
\tightlist
\item
  \textbf{사람 라벨링}: 가장 정확하지만 비용 큼. crowdsourcing, 도메인 전문가, 사내 라벨러.
\item
  \textbf{약지도 학습 (weak supervision)}: 우리 장처럼 휴리스틱 라벨로 \emph{부트스트랩} → 학습된 모델로 \emph{재라벨링} → 사람이 검수. Snorkel 같은 프레임워크.
\item
  \textbf{거대 모델 보조}: GPT-4/Claude 같은 LLM에게 라벨 지시문을 줘서 자동 태깅. 비용은 사람 라벨링보다 훨씬 저렴, 정확도는 도메인에 따라 다름.
\end{enumerate}

이번 합성은 (2)의 가장 단순한 형태입니다 --- 한 번 매칭하고 끝. 실무에선 매칭 결과를 한 번 학습한 모델로 다시 라벨을 매겨 사전을 보강하는 사이클을 돌립니다.

\end{faqBox}

\begin{faqBox}{Q5. (이론) hamming loss는 정확히 무엇을 측정하나요?}

평균 라벨별 오답 비율입니다.

\begin{equation}
\label{eq:ch06-02}
\text{Hamming loss} = \frac{1}{N \cdot K}\sum_{i=1}^{N}\sum_{k=1}^{K} \mathbb{1}[\hat y_{ik} \neq y_{ik}]
\end{equation}

식~\eqref{eq:ch06-02}은 multi-label 평가에서 사용하는 Hamming loss의 정의입니다.

직관: K개 라벨 중 평균적으로 몇 개가 틀렸나. K=5에서 hamming loss 0.1 = 평균 0.5개 라벨 틀림 = 한 샘플의 5개 라벨 중 평균 0.5개 잘못됨.

\textbf{Subset accuracy와의 관계}: subset accuracy는 "5개 라벨 \emph{전부} 맞아야 1점"이라 매우 보수적, hamming은 \emph{부분 점수} 를 줌. 일반적으로 hamming이 더 너그럽게 나옵니다.

\end{faqBox}

\begin{faqBox}{Q6. (실무) 임계값을 라벨마다 다르게 줄 수 있나요?}

네, 자주 합니다. 라벨 빈도가 매우 다를 때 (예: food는 활성률 60\%, location은 10\%) 한 임계값으로는 둘 다 잘 잡기 어렵습니다.

\Needspace{22\baselineskip}
\begin{lstlisting}
from sklearn.metrics import f1_score

# 라벨별 임계값 sweep으로 F1 최대화
best_thresholds = []
for k in range(K):
    best_t, best_f1 = 0.5, 0.0
    for t in np.linspace(0.1, 0.9, 17):
        f1 = f1_score(
            Y_test[:, k],
            (proba_ml[:, k] >= t).astype(int),
            zero_division=0,
        )
        if f1 > best_f1:
            best_t, best_f1 = t, f1
    best_thresholds.append(best_t)

threshold_table = pd.DataFrame({"aspect": ASPECTS, "best_threshold": best_thresholds})
display(threshold_table)
\end{lstlisting}

주의 주의: 위 코드는 \emph{test set} 으로 임계값을 정해 보여준 것 --- 실무에선 \emph{별도 검증 데이터셋} 으로 임계값을 정하고 test에는 적용만 해야 합니다 (안 그러면 test에 누설).

\end{faqBox}

\begin{faqBox}{Q7. (실무) \inlinecode{LogisticRegression(multi\_class="ovr")} (5장)와 \inlinecode{OneVsRestClassifier(LogisticRegression())} (이 장)은 정확히 어떻게 다른가요?}

겉보기엔 같은 OvR 알고리즘인데, 두 도구가 \emph{다른 위치} 에 있는 이유가 있습니다.

\begin{adjustbox}{max width=\textwidth}
\begin{tabular}{@{}lll@{}}
\toprule
비교 항목 & \inlinecode{LogisticRegression(multi\_class="ovr")} (5장) & \inlinecode{OneVsRestClassifier(LogisticRegression())} (6장) \\
\midrule

정체 & 모델 \emph{인자(flag)} & 별도 \emph{wrapper class} (\inlinecode{sklearn.multiclass}) \\
받는 Y 형식 & \textbf{1D만} (multi-class 인덱스 0\textasciitilde K-1) & \textbf{1D + 2D 둘 다} (multi-class 또는 multi-label) \\
내부 저장 & 단일 \inlinecode{coef\_} shape \inlinecode{(K,\ V)} & K개 \emph{별도} estimator를 \inlinecode{.estimators\_} 리스트로, 각자 \inlinecode{coef\_} shape \inlinecode{(1,\ V)} \\
base classifier & \inlinecode{LogisticRegression} 전용 (+ \inlinecode{LogisticRegressionCV}, \inlinecode{LinearSVC} 등 일부 선형 모델) & \textbf{임의 binary 분류기} 가능 (\inlinecode{SVC}, \inlinecode{RandomForestClassifier}, ...) \\
multi-label 지원 & X (1D Y만 받으니 불가능) & OK \\
sklearn 향후 & \inlinecode{multi\_class} 인자 deprecated 진행 중 (1.7+에서 제거 예정) & 표준 클래스로 그대로 유지 \\
\bottomrule
\end{tabular}
\end{adjustbox}

\textbf{핵심 차이}: \inlinecode{multi\_class="ovr"}은 \emph{multi-class 분류 안에서의 옵션} 이고, \inlinecode{OneVsRestClassifier}는 \emph{multi-class와 multi-label 둘 다 처리하는 일반 wrapper} 입니다.

\Needspace{20\baselineskip}
\begin{lstlisting}
# 5장에서: multi-class만 가능 (1D Y)
LogisticRegression(multi_class="ovr").fit(X, y_1d)
# OK 동작
LogisticRegression(multi_class="ovr").fit(X, y_multihot)
# X ValueError (1D 아님)

# 6장에서: multi-class와 multi-label 모두
OneVsRestClassifier(LogisticRegression()).fit(X, y_1d)
# OK multi-class
OneVsRestClassifier(LogisticRegression()).fit(X, y_multihot)
# OK multi-label

# 보너스: base 분류기를 바꾸는 것도 자유
from sklearn.svm import SVC
OneVsRestClassifier(SVC(probability=True)).fit(X, y_multihot)
# OK
\end{lstlisting}

\textbf{언제 무엇을 쓰나}:

\begin{itemize}
\tightlist
\item
  multi-class에서 LogReg의 \emph{내장 OvR} 동작을 잠깐 비교하고 싶다면 \inlinecode{multi\_class="ovr"} (단, 향후 deprecated이라 신규 코드는 권장하지 않음).
\item
  multi-label, 또는 base 분류기를 바꿔야 하거나, K개 binary 모델을 \emph{명시적으로} 다루고 싶다면 \inlinecode{OneVsRestClassifier} (안정적, 미래에도 그대로).
\end{itemize}

이 커리큘럼은 5장에서 \textbf{개념 비교용으로} \inlinecode{multi\_class="ovr"} 을 잠깐 썼고, 6장의 multi-label에서는 \textbf{표준 도구인} \inlinecode{OneVsRestClassifier} 로 정착합니다.
\end{faqBox}



\section{오류 실험 (선택)}

같은 데이터를 multi-class로 잘못 풀면 어떻게 될까요? 측면이 가장 강한 것 \emph{하나} 만 정답으로 골라보고 (argmax) multinomial LogReg를 학습합니다.

\Needspace{18\baselineskip}
\begin{lstlisting}
# 측면이 하나라도 활성된 샘플만 사용
mask_nonempty = Y.sum(axis=1) > 0
y_pseudo_class = Y[mask_nonempty].argmax(axis=1)
# 0-4 사이 정수, "가장 강한 측면"
texts = df["text"][mask_nonempty]

X_pseudo = tfidf.transform(texts)
model_pseudo = LogisticRegression(
    multi_class="multinomial",
    max_iter=1000,
)
model_pseudo.fit(X_pseudo[:4000], y_pseudo_class[:4000])
acc_pseudo = (model_pseudo.predict(X_pseudo[4000:]) == y_pseudo_class[4000:]).mean()
display(pd.DataFrame([{"metric": "pseudo_multiclass_accuracy", "value": acc_pseudo}]))
\end{lstlisting}

힌트: 한 리뷰에 여러 측면이 동시에 있을 때 \emph{하나만} 정답으로 골라 학습하면 정보가 사라집니다. 정답이 임의 선택이라 모델이 어느 라벨을 골라야 할지 모호해지고, multi-label 결과보다 정보가 적은 모델이 됩니다.



\section{Phase 0 마무리 --- sklearn vs HuggingFace 미리보기}

이 장으로 scikit-learn 단계가 끝납니다. 다음 장(7장)부터 등장하는 \inlinecode{transformers} (Hugging Face)는 \emph{loss를 최소화한다} 는 목적은 같지만 \textbf{그 방식과 철학} 이 다릅니다. 큰 그림을 미리 잡아두면 Phase 1의 학습 코드가 익숙하게 이해할 수 있습니다.

\subsection{핵심 차이 한 문장}

\begin{quote}
\textbf{sklearn은 \emph{수학 문제를 풀어준다}.}
\textbf{HuggingFace는 \emph{수학 문제를 푸는 과정} 을 우리가 통제한다.}
\end{quote}

이 장에서 쓴 \inlinecode{LogisticRegression(max\_iter=1000)} 은 lbfgs solver가 알아서 BCE를 최소화하는 가중치를 찾아 돌려줬습니다 --- 우리는 \inlinecode{fit()} 한 줄과 결과만 봤어요.

HF의 \inlinecode{Trainer}는 학습 \emph{과정} 을 명시합니다 --- 학습률, 배치 크기, 에폭 수, 스케줄러, 평가 빈도. loss는 매 step마다 계산되어 backprop으로 가중치를 \emph{조금씩} 옮깁니다. 같은 데이터로 학습해도 random seed가 바뀌면 결과가 미세하게 달라지는 이유.

\subsection{한 표로 정리}

\begin{adjustbox}{max width=\textwidth}
\begin{tabular}{@{}lll@{}}
\toprule
축 & sklearn (Phase 0) & HuggingFace / PyTorch (Phase 1+) \\
\midrule

\textbf{최적화 방식} & 수렴 보장 solver (lbfgs 등) 한 번 호출 → 전역 최적해 & 미니배치 SGD/Adam --- 학습자가 epoch·step 통제 \\
\textbf{언제 끝나나} & 수렴 기준(\inlinecode{tol}) 도달 시 자동 & 사용자 지정 epoch 수 (멈출 시점 직접 결정) \\
\textbf{결정성} & convex 문제라 같은 입력엔 같은 출력 & non-convex --- random seed·batch 순서에 따라 매번 미세 차이 \\
\textbf{에폭/배치 개념} & 보통 없음 --- 전체 데이터 한 번에 & \textbf{핵심} --- \inlinecode{num\_train\_epochs}, \inlinecode{batch\_size} 명시 필수 \\
\textbf{loss를 직접 보나} & 거의 안 봄 (fit 후 평가만) & 매 step마다 loss 출력 + 곡선 추적 (학습이 불안정해지면 즉시 보임) \\
\textbf{하드웨어} & CPU, 단일 스레드 위주 & \textbf{GPU 필수} (fp16, gradient accumulation 등) \\
\textbf{loss 함수 지정} & 모델 클래스에 내장 (LogReg = log loss) & \inlinecode{problem\_type} 자동 매핑 또는 \inlinecode{compute\_loss} 오버라이드 \\
\textbf{모델 크기} & 수만 \textasciitilde{} 수십만 파라미터 & 사전학습 BERT --- 6천만 \textasciitilde{} 수억 파라미터 \\
\textbf{학습 시간 (Yelp 5,000)} & 1초 미만 & T4 GPU에서 2-5분 \\
\bottomrule
\end{tabular}
\end{adjustbox}

\subsection{코드 형태 미리보기 (9장 BERT 회귀에서 본격 등장)}

\Needspace{24\baselineskip}
\begin{lstlisting}
# Phase 0 (sklearn — 우리가 한 형태)
model = LogisticRegression(max_iter=1000)
model.fit(X, y)   # 한 줄. 수렴 기준까지 알아서 풀어줌.

# Phase 1 이후 (HuggingFace — 다음부터)
from transformers import AutoModelForSequenceClassification, TrainingArguments, Trainer

model = AutoModelForSequenceClassification.from_pretrained(
    "distilbert-base-uncased", num_labels=2,
)
args = TrainingArguments(
    output_dir="./output",
    num_train_epochs=3,             # 학습 반복 횟수
    per_device_train_batch_size=16, # 미니배치 크기
    learning_rate=2e-5,             # 학습률
    fp16=True,                      # T4에서 GPU 효율
    logging_steps=20,               # loss 곡선 출력
)
trainer = Trainer(
    model=model,
    args=args,
    train_dataset=...,
    ...,
)
trainer.train()
# 매 step마다 loss → backward → optimizer step
\end{lstlisting}

각 인자가 무엇을 하는지는 9장에서 본격적으로 확인합니다. 지금은 "fit 호출 한 줄이 수십 개 인자로 펼쳐진다" 는 감만 가지면 됩니다.

\subsection{변하지 않는 것}

Loss 자체 --- \inlinecode{BCEWithLogitsLoss}, \inlinecode{CrossEntropyLoss}, \inlinecode{MSELoss} --- 는 sklearn에서 익힌 그대로 Phase 1+에서도 등장합니다. \textbf{달라지는 건 \emph{어떻게 최소화하느냐} 의 도구뿐}. Phase 0의 직관(BCE 수치 표, OvR fit 분해, softmax 동등성 등)이 Phase 1의 BERT 학습에서도 그대로 유지됩니다.



\begin{previewBox}{미리보기: 다음 장}

\textbf{Phase 1 시작 --- 7장. BERT 첫 만남 (\inlinecode{pipeline})}

\begin{itemize}
\tightlist
\item
  scikit-learn 단계 끝, \textbf{\inlinecode{transformers} 라이브러리} 가 처음 등장
\item
  \inlinecode{pipeline("sentiment-analysis")} 한 줄로 사전학습된 DistilBERT 추론
\item
  \inlinecode{pipeline} 안에서 일어나는 3단계 (tokenizer / model / post-processing) 직접 풀어보기
\item
  \textbf{첫 WordPiece 등장} --- 같은 문장이 TF-IDF와 어떻게 다른 단위로 쪼개지는지 비교
\item
  학습 없음 (추론만) --- \inlinecode{Trainer}는 9장에서 본격 등장
\end{itemize}
\end{previewBox}
